% Roteiro para Obtenção de Parâmetros de Detecção de Automação em APIs de Inteligência Artificial
% Formato ABNT (NBR 14724:2011)

\begin{document}

\title{Roteiro para Obtenção de Parâmetros de Detecção de Automação em APIs de Inteligência Artificial: Uma Abordagem Espectral}
\author{NEXUS Project}
\date{2026}

\maketitle

\begin{abstract}
Este documento apresenta um roteiro metodológico para identificação experimental dos parâmetros utilizados por provedores de APIs de Inteligência Artificial na detecção de tráfego automatizado. Fundamentado na análise espectral de séries temporais, o método proposto permite caracterizar padrões de uso humano e diferenciá-los de padrões de automação, mediante transformada rápida de Fourier (FFT) e métricas de entropia espectral. O roteiro inclui procedimentos de captura de tráfego, extração de features, formulação de hipóteses e validação estatística, com vistas à construção de um modelo preditivo de classificação de intenção.
\end{abstract}

\textbf{Palavras-chave:} análise espectral, detecção de automação, APIs de IA, entropia espectral, séries temporais.

\section{Introdução}

Provedores de APIs de Inteligência Artificial (IA) estabelecem limites de uso (\textit{rate limits}) e proíbem automação em seus Termos de Serviço (ToS). Contudo, a detecção de violações opera em espaço multidimensional de características (\textit{features}), cujos parâmetros não são divulgados. Este roteiro formaliza o procedimento para obtenção experimental desses parâmetros, utilizando análise de frequência sobre séries temporais de requisições.

\section{Referencial Teórico}

\subsection{Ruído 1/f (Pink Noise)}
O comportamento humano em interações com sistemas computacionais segue distribuição de potência $P(f) \propto 1/f^{\alpha}$, onde $\alpha \approx 1$ (ruído rosa). Esta propriedade é observada em diversas atividades humanas, desde batimentos cardíacos até padrões de digitação (Voss \& Clarke, 1975).

\subsection{Espectro de Fourier}
A Transformada Rápida de Fourier (FFT) decompõe uma série temporal $x(t)$ em componentes sinusoidais:
$$X(f) = \int_{-\infty}^{\infty} x(t) \cdot e^{-i2\pi ft} dt$$

Padrões periódicos (bots) produzem picos espectrais; padrões estocásticos (humanos) produzem espectro disperso.

\subsection{Entropia Espectral}
A entropia de Shannon aplicada ao espectro de magnitudes mede a uniformidade da distribuição de energia:
$$H = -\sum_{k=1}^{N} p_k \cdot \log_2(p_k)$$
onde $p_k = |X(k)|^2 / \sum |X(j)|^2$. Alta entropia indica espectro plano (humano); baixa entropia indica concentração de energia em frequências específicas (bot).

\section{Metodologia}

\subsection{Fase 1: Captura de Tráfego}

\subsubsection{1.1 Configuração do Ambiente}
\begin{enumerate}
    \item Instalar proxy MITM (mitmproxy) com CA customizado no cliente;
    \item Configurar variáveis de ambiente: \texttt{HTTPS\_PROXY=http://localhost:8080};
    \item Executar clientes legítimos (Claude Code, ChatGPT Web, Gemini) por período mínimo de 7 dias;
    \item Registrar: timestamp, método HTTP, URL, headers, corpo, status de resposta, latência.
\end{enumerate}

\subsubsection{1.2 Amostragem}
\begin{itemize}
    \item Amostragem contínua durante horários de uso ativo;
    \item Mínimo de 10.000 requisições por perfil de usuário;
    \item Segmentação por: provedor, tipo de tarefa (chat, code generation, vision), horário.
\end{itemize}

\subsection{Fase 2: Extração de Features}

\subsubsection{2.1 Features Temporais}
A partir da série de timestamps $\{t_1, t_2, \ldots, t_n\}$, calcular:
\begin{itemize}
    \item Intervalos: $\Delta t_i = t_{i+1} - t_i$
    \item Coeficiente de Variação: $CV = \sigma_{\Delta t} / \mu_{\Delta t}$
    \item Burstiness: $B = (\sigma - \mu) / (\sigma + \mu)$
\end{itemize}

\subsubsection{2.2 Features de Conteúdo}
\begin{itemize}
    \item Tokens por requisição: $T_i$ (request), $T_o$ (response)
    \item Razão: $R_T = T_o / T_i$
    \item Frequência de tool calls: $F_{tool} = N_{tools} / N_{reqs}$
\end{itemize}

\subsection{Fase 3: Análise Espectral}

\subsubsection{3.1 Transformada de Fourier}
\begin{enumerate}
    \item Construir sinal $x(t) = \{\Delta t_1, \Delta t_2, \ldots, \Delta t_n\}$;
    \item Remover tendência: $x'(t) = x(t) - \bar{x}$;
    \item Aplicar janela de Hanning: $w(n) = 0.5 \cdot (1 - \cos(2\pi n / (N-1)))$;
    \item Calcular FFT: $X(k) = \sum_{n=0}^{N-1} x'(n) \cdot w(n) \cdot e^{-i2\pi kn/N}$;
    \item Obter magnitudes: $|X(k)|$ e frequências $f_k = k \cdot f_s / N$.
\end{enumerate}

\subsubsection{3.2 Métricas Espectrais}
\begin{itemize}
    \item Frequência dominante: $f_{dom} = \arg\max_k |X(k)|$ (excluindo DC);
    \item Entropia espectral normalizada: $H_n = H / \log_2(N)$;
    \item \textit{Spectral flatness}: $SF = \exp(\frac{1}{N}\sum \ln|X(k)|) / (\frac{1}{N}\sum |X(k)|)$;
    \item Inclinação espectral: ajuste linear de $\log|X(f)|$ vs $\log f$, obtendo $\alpha$ em $1/f^{\alpha}$.
\end{itemize}

\subsection{Fase 4: Formulação de Hipóteses}

Com base na literatura e observações preliminares, formular:

\begin{itemize}
    \item[H1] \textbf{Threshold Detection}: O provedor bloqueia quando $N_{reqs}/min > X$, independente de conteúdo;
    \item[H2] \textbf{Degradation Curve}: Antes do ban, a latência cresce exponencialmente: $L(t) = L_0 \cdot e^{\lambda t}$;
    \item[H3] \textbf{Token Budget}: Existe orçamento diário $B$ de tokens por chave, e \textit{throttle} inicia em $0.8B$;
    \item[H4] \textbf{Spectral Signature}: Tráfego humano apresenta $\alpha \in [0.8, 1.2]$ (1/f); bot apresenta $\alpha < 0.5$ ou picos dominantes;
    \item[H5] \textbf{Humanoid Timing}: Intervalos com distribuição log-normal ($\mu = 3.4$, $\sigma = 0.5$) não são filtrados.
\end{itemize}

\subsection{Fase 5: Validação}

\subsubsection{5.1 A/B Testing}
\begin{itemize}
    \item Utilizar $N$ chaves idênticas, cada uma com padrão de uso diferente;
    \item Medir \textit{lifetime} até ban: $L_i$ para cada chave $i$;
    \item Correlacionar $L_i$ com métricas espectrais via regressão linear múltipla.
\end{itemize}

\subsubsection{5.2 Análise de Sobrevivência}
Modelar tempo até ban usando distribuição Weibull:
$$h(t) = \frac{k}{\lambda} \left(\frac{t}{\lambda}\right)^{k-1}$$
onde $k$ (forma) e $\lambda$ (escala) são estimados por máxima verossimilhança.

\subsubsection{5.3 Significância Estatística}
\begin{itemize}
    \item Teste de Mann-Whitney para comparação de distribuições (humano vs bot);
    \item $p$-value $< 0.05$ para rejeição de hipóteses nulas;
    \item Intervalo de confiança de 95\% para parâmetros estimados.
\end{itemize}

\section{Resultados Esperados}

Ao final do roteiro, obter-se-á:

\begin{enumerate}
    \item Conjunto de parâmetros $\{\alpha, H_n, SF, f_{dom}\}$ que caracterizam tráfego humano;
    \item Modelo preditivo de probabilidade de detecção: $P(ban | X) = \sigma(\beta^T X)$;
    \item \textit{Thresholds} ótimos para operação segura: $\theta_{opt} = \arg\max_{\theta} (throughput \cdot P(não\ ban))$;
    \item Estratégia de diluição temporal baseada em síntese de ruído 1/f.
\end{enumerate}

\section{Considerações Éticas}

Este roteiro destina-se exclusivamente a:
\begin{itemize}
    \item Compreensão de mecanismos de detecção para operação dentro dos termos de uso;
    \item Pesquisa acadêmica sobre sistemas de detecção de anomalias;
    \item Desenvolvimento de ferramentas de diagnóstico de tráfego próprio.
\end{itemize}

Não se destina a:
\begin{itemize}
    \item Evasão deliberada de sistemas anti-fraude;
    \item Compartilhamento não autorizado de credenciais;
    \item Violação de Termos de Serviço de terceiros.
\end{itemize}

\section{Referências}

VOSS, R. F.; CLARKE, J. 1/f noise in music: Music from 1/f noise. \textbf{Journal of the Acoustical Society of America}, v. 63, n. 1, p. 258-263, 1978.

SHANNON, C. E. A mathematical theory of communication. \textbf{Bell System Technical Journal}, v. 27, n. 3, p. 379-423, 1948.

MANN, H. B.; WHITNEY, D. R. On a test of whether one of two random variables is stochastically larger than the other. \textbf{The Annals of Mathematical Statistics}, v. 18, n. 1, p. 50-60, 1947.

NBR 14724:2011. \textbf{Informação e documentação — Trabalhos acadêmicos — Apresentação}. ABNT, 2011.

\end{document}
